% Section 1 (Introduction) draft for "Beyond Web5: A Substrate Map"
% Locked 2026-04-29; operator-revised LaTeX from canon
% Source: project_neurips_paper_pivot_substrate_map_beyond_web5.md
% Markdown sibling: airlock-labs-site/papers/beyond-web5-substrate-map.md

\section{Introduction}

The last decade of web and AI infrastructure has produced a familiar ladder:
Web1 (read), Web2 (read--write), Web3 (own), and Web4 (inhabit, via
substrate-resident agents) \cite{bernerslee1989,oconnor2023web4}. Recent proposals
branded as ``Web5'' argue for an additional layer centered on decentralized
identity and personal-server federation \cite{tbd_web5,solid_pods,ssi_survey}, but
largely stop at architectural slogans rather than concrete, measurable systems.
In this work we take a complementary approach: instead of debating the exact
definition of Web5, we implement and instrument an operator-resident substrate
that sits \emph{beyond} the conventional ladder, treating Web1--4 as input
surfaces to a single local control shell.
The resulting system combines seven primitives---identity, logging, behavioral
scoring, inheritance, incentives, and federation---into a working ``operator OS''
whose behavior we can probe empirically using a 17-persona benchmark
(ConstellationBench) and a fully wired state service with 4/4 health, loading,
persistence, and warm-start checks passing.

\subsection{From Web5 slogans to operator-resident substrates}

Most Web5 discussions to date focus on protocol proposals: decentralized
identifiers (DIDs) and verifiable credentials \cite{w3c_did_core}, personal data
stores and Solid pods \cite{solid_pods}, or high-level visions of ``the web of
identity'' \cite{tbd_web5}. These lines of work are important, but they leave
open a concrete systems question: \emph{what primitives are actually required
to operate a real, operator-resident substrate that sits above existing web
layers and interacts with them in practice?}

We address this question empirically. Rather than specifying a new protocol
and leaving its deployment as future work, we start from a running system used
for day-to-day development, then factor out the primitives that proved necessary
to make it observable, auditable, and controllable by a human operator.
Concretely, our implementation exposes seven reusable primitives:
\begin{enumerate}
  \item phone-anchored decentralized identifiers (Skeet);
  \item a biometric-gated local control shell that federates identities
        (Spacebar v0);
  \item signed bivector closures for event logging (BCODE);
  \item behavioral-fidelity scoring via DECF and a 17-persona dataset
        (ConstellationBench);
  \item agent inheritance patterns for specialized assistants (OctoAgent);
  \item a token-economy layer for incentives (\$MOTION with ACODE/BCODE split);
  \item matrix-style federation transport (Matrix / bitchat).
\end{enumerate}
Each primitive is implemented in a real codebase, exported as an open artifact,
and evaluated in the context of the running substrate.

\subsection{ConstellationBench as an empirical lens}

ConstellationBench is a 17-persona benchmark designed to measure
behavioral fidelity of large language models along multiple DECF axes:
directionality, explicitness, coherence, and fidelity \cite{constellationbench}.
Rather than being the sole focus of this paper, ConstellationBench plays the
role of an \emph{empirical lens} on the substrate: we use it to probe how
different primitives behave when populated with real personas and real models.

First, we use DECF scores across personas to quantify how consistently
frontier models reproduce intended behavioral profiles, providing a concrete
notion of ``behavioral fidelity'' that goes beyond aggregate task accuracy.
Second, we instrument the substrate so that every evaluation run produces
signed BCODE closures and serialized context, allowing us to study how
inheritance patterns (OctoAgent) and logging interact in practice.
The Phase-0 wiring---a state service, persona loader, session persistence, and
warm-start behavior, with 4/4 canonical wires green and 25/26 tests passing---serves
as our reproducibility receipt: it specifies exactly which components must be
operational for other groups to replicate our experiments or to plug in
alternative primitives.

\subsection{Contributions}

Our contributions are threefold:
\begin{itemize}
  \item \textbf{Substrate map.} We identify seven primitives that, in combination,
        are sufficient to operate a working operator-resident substrate above
        existing web layers, and we provide an architectural map showing how
        they compose into a single local control shell.
  \item \textbf{Empirical lens.} We use ConstellationBench as a behavioral-fidelity
        lens on this substrate, reporting results for multiple frontier models
        across 17 personas and analyzing how inheritance and logging primitives
        behave under stress.
  \item \textbf{Open artifacts.} We release the ConstellationBench dataset,
        evaluation runtime, and core substrate components (ACODE-free) so that
        other operator-resident systems can reuse individual primitives or
        benchmark alternative designs against the same empirical lens.
\end{itemize}
